%% kapa: exam=2 round=20 year=2009 period=3 subject=compensation_law subject_ko=감정평가 및 보상법규 seq=1 number=1 points=45
%% source: Q-Net 감정평가사 기출문제 (cst00309) artlSeq=5024180 file=extracted/1008/5024180/2009년 제20회 감정평가사 2차시험 기출문제/3교시.hwp
%% conversion: kordoc 4.24.51 -> tools/gen2.py (auto); review=proofread-sonnet; flags=none
\begin{longquestion}{【문제 1】}
A시는 도시개발사업을 하면서 주거를 상실하는 거주자에 대한 이주대책을 수립하였다. 이주대책의 주요내용은 다음과 같다. 이를 근거로 다음 물음에 답하시오. (45점) \par
\begin{framedtext}
○ 기준일 이전부터 사업구역 내 자기 토지상 주택을 소유하고 협의계약 체결일까지 당해 주택에 계속 거주한 자가 보상에 합의하고 자진 이주한 경우 사업구역 내 분양아파트를 공급한다. \par
○ 분양아파트를 공급받지 않은 이주자에게는 이주정착금을 지급한다. \par
○ 무허가건축물 대장에 등록된 건축물 소유자는 이주대책에서 제외한다.
\end{framedtext}
\medskip\noindent (1) 이주대책의 이론적 및 헌법적 근거를 설명하시오. (5점) \par
\medskip\noindent (2) 주택 소유자 甲이 보상에 합의하고 자진 이주하지 아니한 경우에도 이주대책에 의한 분양아파트의 공급 혹은 이주정착금의 지급을 요구할 수 있는지의 여부를 검토하시오. (20점) \par
\medskip\noindent (3) 무허가건축물대장에 등록되지 않은 건축물 소유자 乙이 당해 건축물이 무허가건축물이라는 이유로 이주대책에서 제외된 경우에 권리구제를 위하여 다툴 수 있는 근거와 소송방법에 관하여 검토하시오. (20점) \par
\end{longquestion}
